\scanpage{29}
\unsure{noté} $C^*K$, défini par les générateurs
\[
c_{\deleted{\unsure{i}}}^{i}(x)\qquad i\geq1
\]
et les relations
\[
\left\{\begin{array}{ll}
c^i(1)=0 & \text{si }i\geq1,\\
c^n(x+y)=\displaystyle\sum_{i+j=n}c^i(x)c^j(y)&\text{si }x,y\in K,\ n\geq1\\
c^n(x*y)=Q_n(c^1(x),\text{---},c^n(x);c^1(y),\text{---},c^n(y))&\text{pour }x,y\in I(K),\ n\geq1\\
c^n(E_s)=0&\text{pour }n>n_s.
\end{array}\right.
\]
\margin{\unsure{par conséquence des autres relations}}

Noter que l'on a pour tout $E_s$
\[
\gamma^n(E_s-n_s)=0\qquad\text{pour }n>n_s
\]
[en vertu de $\lambda^i(E_s)=0$ pour $i>n_s$] donc les \uline{classes de Chern \unsure{naturelles}} satisfont aux conditions de classes de Chern \deleted{géométriques}\inserted{admissibles} pour les $E_s$, on a donc un hom can
\[
\varphi_K:CK\longrightarrow GK
\]
\[
c^i(x)\rightsquigarrow\gamma^i(x-\varepsilon(x))\mod K^{(i+1)}
\]
évidemment \uline{surjectif} [N.B. plus la famille des $E_s$ est grande, plus on impose des conditions, donc plus $CK$ devient petit, plus on \unsure{peut espérer} demander que $\varphi_K$ soit un \unsure{isom}].

\margin{À voir: la filtration de $K$ définie par les $P(\gamma^i(x-\varepsilon(x)))$ est aussi celle définie par les $P(\gamma^i(E_s-\varepsilon_s))$, [si les $E_s$ engendrent] ???}

\margin{N.B. Il n'y a pas lieu en général, (pour les applications) de \unsure{remplacer} la famille des $E_s$ par l'ens de \unsure{ts les éléments} de $K$ qui sont ``géométriques'', \unsure{car il pourra alors y avoir des} ``bonnes'' applications $\unsure{c_i|c^i}$ telles que $\unsure{c^i\ne0}$ \unsure{pour $i>\varepsilon$}!}

\deleted{D'ailleurs, \unsure{on a tenté de montrer}}
\margin{28}
